\chapter{TỔNG KẾT}\label{chap4_intro}
\vspace{-40pt}

\section{Kết quả đạt được}

\subsection{Hoàn thành các chức năng chính của hệ thống}

Các chức năng đã triển khai bao gồm:
\begin{itemize}
    \item Đăng ký, đăng nhập, đăng xuất bằng tài khoản của hệ thống; hỗ trợ đăng nhập và liên kết tài khoản bằng Google.
    \item Xem trang chủ, danh mục, danh sách sản phẩm, tìm kiếm, lọc và xem thông tin chi tiết sản phẩm.
    \item Quản lý giỏ hàng theo tài khoản khách hàng, thêm hoặc xóa sản phẩm, thay đổi số lượng và tính lại tổng tiền.
    \item Quản lý hồ sơ khách hàng và sổ địa chỉ giao hàng; cho phép người dùng cập nhật thông tin cá nhân, mật khẩu và địa chỉ.
    \item Tạo đơn hàng thanh toán khi nhận hàng (COD), chọn địa chỉ, kiểm tra tồn kho, tạm trừ số lượng khi đặt hàng và xóa các dòng giỏ hàng đã chuyển thành đơn.
    \item Xem danh sách đơn hàng, xem chi tiết đơn và hủy đơn ở trạng thái cho phép; hoàn lại tồn kho khi đơn được hủy hoặc hàng hoàn về kho theo quy tắc xử lý.
    \item Tích hợp SePay ở phương thức thanh toán trực tuyến thông qua sandbox, tạo yêu cầu thanh toán và xử lý kết quả thành công, thất bại, hủy giao dịch cùng thông báo IPN.
    \item Quản lý sản phẩm, danh mục, người dùng, nhân viên, tồn kho, khuyến mãi, voucher, đơn hàng và thống kê kinh doanh trong khu vực quản trị.
    \item Quản lý trạng thái đơn hàng theo quy trình từ Chờ xác nhận, Đang chuẩn bị hàng, Đang giao hàng đến Hoàn thành hoặc Giao hàng thất bại.
    \item Ghi nhận lịch sử thay đổi tồn kho, gồm loại thao tác, sản phẩm, số lượng thay đổi, người thực hiện và thời gian.
\end{itemize}

\subsection{Đáp ứng các yêu cầu về dữ liệu và nghiệp vụ}

Mã nguồn đã sử dụng prepared statements của PDO để truy vấn dữ liệu dựa trên đầu vào của người dùng. Các thao tác thay đổi dữ liệu quan trọng được xử lý trong transaction, đặc biệt là đăng ký tài khoản, cập nhật giỏ hàng, tạo đơn hàng, hủy đơn và cập nhật tồn kho. Khi một giao dịch gặp lỗi, hệ thống rollback để tránh lưu một phần dữ liệu hoặc làm thay đổi tổng thể sản phẩm không nhất quán.

Mật khẩu người dùng được lưu bằng cơ chế \texttt{password\_hash()} và kiểm tra bằng \texttt{password\_verify()}. Các biểu mẫu thay đổi dữ liệu sử dụng mã CSRF. Controller kiểm tra phiên đăng nhập, vai trò và trạng thái tài khoản trước khi cho truy cập các chức năng giỏ hàng, đặt hàng, tài khoản hoặc khu vực quản trị.

Cơ sở dữ liệu MySQL lưu các nhóm dữ liệu chính gồm người dùng, tài khoản, danh mục, sản phẩm, giỏ hàng, đơn hàng, địa chỉ nhận hàng, voucher, khuyến mãi, lịch sử kho và thông tin thanh toán. Các khóa chính, khóa ngoại và ràng buộc được sử dụng để liên kết dữ liệu và kiểm soát quan hệ nghiệp vụ.

Hệ thống đã áp dụng một số quy tắc kinh doanh quan trọng:
\begin{itemize}
    \item Không cho thêm sản phẩm ngừng kinh doanh hoặc vượt quá số lượng tồn kho.
    \item Tính lại giá khuyến mãi và voucher ở phía máy chủ; không sử dụng số tiền do trình duyệt gửi lên làm cơ sở tạo đơn.
    \item Chỉ ghi nhận lượt sử dụng voucher khi đơn hàng được tạo thành công.
    \item Chỉ khách hàng sở hữu đơn mới được xem và hủy đơn; chỉ hủy được đơn ở trạng thái Chờ xác nhận.
    \item Doanh thu thống kê chỉ được tính từ đơn hàng ở trạng thái Hoàn thành.
    \item Đơn hàng giữ lịch sử tổng tiền, địa chỉ nhận hàng và giá trị voucher tại thời điểm đặt hàng.
\end{itemize}

\subsection{Đánh giá bước đầu}

Các thành phần chức năng đã được tích hợp thành luồng sử dụng liên tục thay vì xây dựng các trang giao diện độc lập. Người dùng có thể xem sản phẩm, quản lý giỏ hàng, đặt hàng, theo dõi đơn và thực hiện các thao tác liên quan đến tài khoản. Phía quản trị có thể xử lý đơn hàng và quản lý dữ liệu kinh doanh trong phạm vi website đã xây dựng.

Các kết quả kiểm tra hiện tại chủ yếu dựa trên đối chiếu mã nguồn, cấu trúc route, controller, model và giao diện. Các luồng xử lý thành công và thất bại đã được tổ chức trong controller, có kiểm tra dữ liệu đầu vào.

\section{Hạn chế và hướng phát triển}

\subsection{Hạn chế hiện tại}
Vẫn còn một số hạn chế mặc dù nhóm đã cố gắng hoàn thiện đồ án, nhưng do kiến thức và kinh nghiệm thực tế của sinh viên còn hạn chế nên đề tài vẫn còn một số khuyết điểm. Chưa tiến hành kiểm thử trực tiếp trên môi trường Apache và cơ sở dữ liệu MySQL với dữ liệu thực, nên kết quả của đồ án cần được xác nhận lại bằng các bài kiểm thử chạy hệ thống thực tế.

\subsection{Hướng phát triển}

% Trong thời gian tiếp theo, đồ án nên bổ sung bộ dữ liệu kiểm thử và xây dựng test case cho các luồng đăng ký, đăng nhập, giỏ hàng, đặt hàng, hủy đơn, thanh toán, quản lý kho và phân quyền. Đồng thời, cần tiến hành chạy thử thực tế trên XAMPP, kiểm tra dữ liệu sau khi tạo, sửa, hủy hoặc hoàn đơn và ghi lại kết quả dưới dạng log hoặc màn hình minh họa.

Hệ thống có thể được hoàn thiện bằng các hướng phát triển sau:
\begin{itemize}
    \item Bổ sung bộ lọc nâng cao, tìm kiếm toàn văn hoặc đánh chỉ mục cơ sở dữ liệu để cải thiện hiệu năng khi số lượng sản phẩm tăng.
    \item Hoàn thiện giao diện trên các kích thước màn hình và trình duyệt phổ biến.
    \item Hoàn thiện cấu hình, nhật ký và cơ chế giám sát thanh toán trực tuyến, đồng thời kiểm tra bảo mật callback/IPN trong môi trường production.
    \item Bổ sung tự động hóa triển khai, sao lưu cơ sở dữ liệu và hướng dẫn cài đặt, import schema và migration.
    \item Mở rộng báo cáo thống kê theo doanh thu, chi phí, lợi nhuận, tình trạng kho, hình thức thanh toán và từng khoảng thời gian.
    \item Cải thiện xử lý ảnh sản phẩm, giới hạn dung lượng tải lên và quản lý hình ảnh tập trung nếu chức năng upload được phát triển.
\end{itemize}

Tổng kết lại, đồ án đã đạt được mục tiêu ứng dụng các kiến thức đã học từ các môn học để xây dựng website thương mại điện tử có luồng mua hàng và khu vực quản trị cơ bản, đồng thời kết nối nghiệp vụ với cơ sở dữ liệu MySQL. Kết quả còn cần được hoàn thiện bằng kiểm thử chạy thực tế, xử lý dữ liệu kiểm thử và bổ sung các chức năng mở rộng để hệ thống có thể đáp ứng tốt hơn nhu cầu sử dụng thực tế.
